%! TEX root = ../main.tex

\usepackage{tocloft}

% 0. CHO PHÉP MỤC LỤC HIỂN THỊ ĐẾN HEADING 4 (\paragraph)
\setcounter{tocdepth}{4}

% 1. TIÊU ĐỀ MỤC LỤC
\renewcommand{\cfttoctitlefont}{\hfill\LARGE\bfseries}
\renewcommand{\contentsname}{MỤC LỤC}
\renewcommand{\cftaftertoctitle}{\hfill}

% --- GIẢM KHOẢNG TRỐNG PHÍA TRÊN HEADING 1 (SECTION) ---
% Mặc định là khoảng 1.0em + 0.6pt. Bạn có thể hạ xuống 0.2em hoặc 0pt tùy ý:
\setlength{\cftbeforesecskip}{0.2em}

% 2. DẤU NỐI & KHOẢNG CÁCH DẤU CHẤM
\renewcommand{\cftsecleader}{\cftdotfill{\cftdotsep}}
\renewcommand{\cftdotsep}{2} 

% Hiển thị dấu chấm nối sang số trang cho cả Section (Chương)
\renewcommand{\cftsecpagefont}{\normalfont}
\renewcommand{\cftsecpresnum}{}

% 3. THỤT LỀ TRONG MỤC LỤC
% Syntax: \cftsetindents{entry}{indent}{numwidth}
\cftsetindents{section}{0em}{1.5em}        % Dành không gian cho chữ "CHƯƠNG I:"
\cftsetindents{subsection}{1.5em}{2em}    % Thụt lề cấp 2
\cftsetindents{subsubsection}{3.5em}{2.5em}% Thụt lề cấp 3
\cftsetindents{paragraph}{5.5em}{1.5em}   % Thụt lề cấp 4 (dành cho nhãn a), b)...)



% --- 1. ĐỔI TÊN TIÊU ĐỀ TIẾNG VIỆT ---
\renewcommand{\listfigurename}{DANH MỤC HÌNH}
\renewcommand{\listtablename}{DANH MỤC BẢNG}

% --- 2. ĐỊNH DẠNG TIÊU ĐỀ (Căn giữa, cỡ chữ LARGE, in đậm) ---
\renewcommand{\cftloftitlefont}{\hfill\LARGE\bfseries}
\renewcommand{\cftafterloftitle}{\hfill}

\renewcommand{\cftlottitlefont}{\hfill\LARGE\bfseries}
\renewcommand{\cftafterlottitle}{\hfill}

% --- 3. DẤU NỐI VÀ KHOẢNG CÁCH DẤU CHẤM (Đồng bộ với TOC) ---
\renewcommand{\cftfigleader}{\cftdotfill{\cftdotsep}}
\renewcommand{\cfttableader}{\cftdotfill{\cftdotsep}}

% --- 4. THỤT LỀ TRONG DANH MỤC HÌNH VÀ BẢNG ---
% Cú pháp: \cftsetindents{loại}{lùi_lề_trái}{độ_rộng_dành_cho_nhãn}
% Độ rộng nhãn 4.0em là đủ rộng cho các số dạng "Hình I.1:" hoặc "Bảng II.1:"

% Dành cho danh sách Hình (figure)
\cftsetindents{figure}{0em}{4.0em}

% Dành cho danh sách Bảng (table)
\cftsetindents{table}{0em}{4.0em}

% --- 5. CẤU HÌNH CHO PDF (Hyperref) ---
\hypersetup{
    colorlinks=true,      % Bật màu sắc cho link thay vì đóng khung chữ nhật
    linkcolor=black,      % Màu của các liên kết trong mục lục (để black nếu muốn in ấn chuyên nghiệp, hoặc blue/navy)
    citecolor=black,      % Màu của liên kết trích dẫn tài liệu tham khảo
    urlcolor=blue,        % Màu của liên kết website (URL)
    bookmarksnumbered=true% Hiển thị số thứ tự chương/mục trong thanh Bookmark của phần mềm đọc PDF
}